\documentclass[11pt,a4paper]{article}

\usepackage[margin=1in]{geometry}
\usepackage{amsmath,amssymb,amsthm,mathtools}
\usepackage{booktabs}
\usepackage{enumitem}
\usepackage{hyperref}
\usepackage{xcolor}

\title{\textbf{COSA: Conic Active-Set Algorithm}\\
\large Full Project Plan for a Conic Active-Set Solver}
\author{}
\date{September 2026}

\begin{document}

\maketitle

\begin{abstract}
This project develops COSA, a Conic Active-Set Algorithm for
second-order cone programs (SOCPs), with mean--standard-deviation
portfolio optimization as the primary application. The central idea
is to extend the classical active-set philosophy from polyhedral
constraints to second-order cones. The portfolio problem has a
polyhedral set of trading and exposure constraints, while standard
deviation enters naturally through a second-order cone. The project
will develop the mathematical foundations, numerical algorithm,
implementation, and computational validation of this approach.

The main research question is whether conic optimization problems can
be solved effectively by explicitly identifying and maintaining the
active geometry of the solution, rather than relying exclusively on
interior-point methods. Particular emphasis will be placed on exact
cone geometry, working-set updates, conic KKT conditions, numerical
linear algebra, and warm starts for sequences of related portfolio
problems.
\end{abstract}

\tableofcontents

\newpage

% ============================================================
\section{Project Overview}
% ============================================================

\subsection{Motivation}

Classical active-set methods are particularly attractive for
optimization problems with a polyhedral feasible region. They
maintain a working set of constraints that are believed to be active
at the solution and repeatedly solve equality-constrained subproblems
while adding and removing constraints.

Portfolio optimization provides a natural setting in which this
philosophy is highly relevant. Typical portfolio constraints are
linear:
\[
\mathbf{1}^{\top}x=1,
\qquad
x\geq 0,
\qquad
\ell\leq x\leq u,
\qquad
Ax\leq b.
\]

However, the natural risk measure of interest is often standard
deviation rather than variance:
\[
\sigma(x)
=
\sqrt{x^{\top}\Sigma x}.
\]

Standard deviation has the same physical units as expected return.
Consequently, a risk-adjusted objective of the form
\[
\mu^{\top}x-\lambda\sigma(x)
\]
has a direct interpretation in terms of return and risk.

The resulting optimization problem is convex but not a quadratic
program. The purpose of COSA is to investigate whether its special
second-order cone structure can be incorporated directly into an
active-set framework.

\subsection{Primary Research Question}

The central question of the project is:

\begin{quote}
Can the classical active-set philosophy for polyhedral optimization
be extended to second-order cone constraints in a way that is
mathematically principled, numerically robust, and useful for
structured portfolio optimization?
\end{quote}

The intended answer should include both a mathematical algorithm and
a working implementation.

\subsection{Guiding Principle}

The project is guided by the following principle:
\[
\boxed{
\text{Do for cones what classical active-set methods do for
polyhedra.}
}
\]

The objective is not simply to solve an SOCP using an existing
general-purpose method. The goal is to understand and exploit the
geometry of the active cone directly.

% ============================================================
\section{Mathematical Problem}
% ============================================================

\subsection{Mean--Standard-Deviation Portfolio Problem}

Let
\[
x\in\mathbb{R}^n
\]
denote portfolio weights, let
\[
\mu\in\mathbb{R}^n
\]
denote expected returns, and let
\[
\Sigma\succeq0
\]
denote the covariance matrix.

The portfolio expected return is
\[
r(x)=\mu^\top x,
\]
and portfolio standard deviation is
\[
\sigma(x)
=
\sqrt{x^\top\Sigma x}.
\]

We consider
\begin{equation}
\begin{aligned}
\max_x\quad&
\mu^\top x-\lambda\sqrt{x^\top\Sigma x}\\
\text{s.t.}\quad&
Ax\leq b,\\
&Ex=d,
\end{aligned}
\tag{1}
\end{equation}
where $\lambda>0$ controls the trade-off between expected return and
risk.

Equivalently, we minimize
\[
-\mu^\top x+\lambda\sqrt{x^\top\Sigma x}.
\]

\subsection{Second-Order Cone Representation}

Choose a factorization
\[
\Sigma=L^\top L.
\]

Then
\[
\sqrt{x^\top\Sigma x}
=
\|Lx\|_2.
\]

Introduce an auxiliary variable $t$. The problem becomes
\begin{equation}
\begin{aligned}
\min_{x,t}\quad&
-\mu^\top x+\lambda t\\
\text{s.t.}\quad&
Ax\leq b,\\
&Ex=d,\\
&\|Lx\|_2\leq t.
\end{aligned}
\tag{2}
\end{equation}

Define the second-order cone
\[
\mathcal Q^{m+1}
=
\{(t,y)\in\mathbb{R}^{m+1}:\|y\|_2\leq t\}.
\]

The final constraint can therefore be written as
\[
(t,Lx)\in\mathcal Q^{m+1}.
\]

Thus the complete problem is an SOCP with a polyhedral component
and a second-order cone component.

\subsection{Why Standard Deviation}

Variance gives the quadratic objective
\[
\mu^\top x-\gamma x^\top\Sigma x,
\]
which is particularly convenient computationally.

However, variance has units of return squared, while expected return
has units of return. Standard deviation has the same units as
expected return:
\[
[\sigma]=[\mu].
\]

Therefore
\[
\mu^\top x-\lambda\sigma(x)
\]
is directly interpretable as a return-minus-risk objective.

The project deliberately accepts the additional conic structure in
order to preserve this interpretation.

% ============================================================
\section{COSA Concept}
% ============================================================

\subsection{Classical Active-Set Philosophy}

For a polyhedral problem
\[
\min_x f(x)
\quad\text{s.t.}\quad
Ax\leq b,\quad Ex=d,
\]
the active set at $x_k$ is
\[
\mathcal A_k
=
\{i:a_i^\top x_k=b_i\}.
\]

The working set imposes
\[
a_i^\top p=0,
\qquad i\in\mathcal A_k,
\]
and
\[
Ep=0.
\]

A direction is then computed within the tangent space of the
currently active constraints.

\subsection{Extension to a Cone}

For COSA, the working set contains:

\begin{enumerate}
    \item active linear inequalities;
    \item equality constraints;
    \item the currently active geometry of the second-order cone.
\end{enumerate}

Let
\[
z=(t,Lx).
\]

The SOC constraint is
\[
z\in\mathcal Q.
\]

At a nonzero boundary point,
\[
t=\|Lx\|_2,
\]
define
\[
u=\frac{Lx}{\|Lx\|_2}.
\]

The tangent condition is
\begin{equation}
\tau-u^\top Lp=0.
\tag{3}
\end{equation}

This equation is the local analogue of an active linear constraint.

\subsection{Important Geometric Qualification}

A second-order cone is not simply a nonlinear scalar inequality.
Its geometry is naturally described through tangent cones, normal
cones, and faces.

Therefore COSA will initially use the tangent representation to
construct directions, but the final algorithm will be formulated in
terms of the primal-dual conic KKT conditions.

This distinction is important. The project should not merely become
an SQP method with an SOC constraint. The objective is a genuine
conic active-set method.

% ============================================================
\section{Algorithmic Architecture}
% ============================================================

\subsection{Iteration Structure}

The fundamental COSA iteration is

\[
\boxed{
\text{working set}
\rightarrow
\text{direction}
\rightarrow
\text{feasible step}
\rightarrow
\text{multiplier test}
\rightarrow
\text{working-set update}.
}
\]

At iteration $k$:

\begin{enumerate}
    \item Evaluate primal feasibility.
    \item Identify active linear constraints.
    \item Determine the SOC status.
    \item Construct the current working-set equations.
    \item Solve the direction problem.
    \item Determine the maximum feasible step.
    \item Update the primal variables.
    \item Compute multipliers.
    \item Add or remove constraints.
    \item Check the conic KKT conditions.
\end{enumerate}

\subsection{Direction Problem}

Let
\[
f(x,t)=-\mu^\top x+\lambda t.
\]

For a candidate direction $(p,\tau)$, a regularized direction problem
can be written as
\begin{equation}
\begin{aligned}
\min_{p,\tau}\quad&
\nabla f^\top
\begin{bmatrix}
p\\
\tau
\end{bmatrix}
+
\frac12
\begin{bmatrix}
p\\
\tau
\end{bmatrix}^{\!\top}
H
\begin{bmatrix}
p\\
\tau
\end{bmatrix}\\
\text{s.t.}\quad&
a_i^\top p=0,
\qquad i\in\mathcal A_k,\\
&Ep=0,\\
&\tau-u_k^\top Lp=0
\qquad\text{if the SOC is active.}
\end{aligned}
\tag{4}
\end{equation}

The initial implementation may use
\[
H=\rho I,
\qquad \rho>0,
\]
to obtain a well-defined direction.

\subsection{KKT System}

The direction problem leads to a linear KKT system of the form
\[
\begin{bmatrix}
H&W_k^\top\\
W_k&0
\end{bmatrix}
\begin{bmatrix}
d_k\\
\nu_k
\end{bmatrix}
=
-
\begin{bmatrix}
g_k\\
0
\end{bmatrix},
\]
where
\[
d_k=
\begin{bmatrix}
p_k\\
\tau_k
\end{bmatrix}
\]
and $W_k$ contains the current working-set equations.

This creates a direct connection with classical active-set QP
implementations.

% ============================================================
\section{Exact SOC Step Calculation}
% ============================================================

\subsection{Cone Feasibility Along a Direction}

Given a feasible point $(x,t)$ and a direction $(p,\tau)$, we require
\[
\|L(x+\alpha p)\|_2
\leq
t+\alpha\tau.
\tag{5}
\]

Define
\[
r=Lx,
\qquad
q=Lp.
\]

Then
\[
\|r+\alpha q\|_2
\leq
t+\alpha\tau.
\]

Provided the right-hand side is nonnegative, squaring gives
\begin{equation}
\|q\|_2^2\alpha^2
+
2(r^\top q-t\tau)\alpha
+
\|r\|_2^2-t^2
\leq0.
\tag{6}
\end{equation}

Thus the SOC feasibility condition reduces to a scalar quadratic
inequality.

\subsection{Maximum Feasible Step}

The COSA implementation will compute the admissible interval in
$\alpha$ from (6), together with the step bounds induced by the
linear inequalities.

For each inactive linear inequality,
\[
a_i^\top(x+\alpha p)\leq b_i
\]
gives the standard bound
\[
\alpha
\leq
\frac{b_i-a_i^\top x}{a_i^\top p}
\]
whenever $a_i^\top p>0$.

The feasible step is therefore determined by the intersection of:

\begin{itemize}
    \item the linear feasible interval;
    \item the SOC feasible interval;
    \item any explicit step bounds.
\end{itemize}

This is the conic analogue of the classical active-set ratio test.

% ============================================================
\section{Conic KKT Conditions}
% ============================================================

\subsection{Primal Problem}

The primal SOCP is
\begin{equation}
\begin{aligned}
\min_{x,t}\quad&
-\mu^\top x+\lambda t\\
\text{s.t.}\quad&
Ax\leq b,\\
&Ex=d,\\
&(t,Lx)\in\mathcal Q.
\end{aligned}
\tag{7}
\end{equation}

Introduce multipliers $y\geq0$ for the linear inequalities and
$\nu$ for the equalities. Let $w$ denote the dual variable associated
with the SOC.

\subsection{Stationarity}

The conic stationarity equations have the form
\[
-\mu+A^\top y+E^\top\nu+L^\top w=0,
\]
together with the stationarity condition associated with $t$.

The exact sign convention will be fixed in the implementation and
must be used consistently throughout the derivation and code.

\subsection{Primal and Dual Cone Feasibility}

The primal SOC condition is
\[
(t,Lx)\in\mathcal Q.
\]

The dual SOC variable must satisfy
\[
w_{\mathrm{soc}}\in\mathcal Q,
\]
under the chosen standard representation.

\subsection{Complementarity}

For the linear constraints,
\[
y_i(a_i^\top x-b_i)=0.
\]

For the SOC, complementarity is expressed through the cone
complementarity relation between the primal and dual cone variables.

The final implementation will use residuals for:

\begin{itemize}
    \item primal feasibility;
    \item dual feasibility;
    \item stationarity;
    \item linear complementarity;
    \item SOC complementarity.
\end{itemize}

These residuals define the primary termination criterion.

% ============================================================
\section{Active-Set Logic}
% ============================================================

\subsection{Linear Constraint Activation}

If an inactive inequality reaches its boundary,
\[
a_i^\top x=b_i,
\]
the constraint is added to the working set.

The corresponding multiplier is subsequently computed from the
KKT system.

\subsection{Linear Constraint Deactivation}

If an active inequality has a multiplier violating the required sign,
the constraint is a candidate for removal.

The implementation will initially follow the classical rule of
removing the most strongly violating multiplier, subject to
numerical tolerances.

\subsection{SOC Activation}

The SOC becomes geometrically active when
\[
t-\|Lx\|_2
\]
is sufficiently small.

A tolerance such as
\[
t-\|Lx\|_2\leq\varepsilon_{\mathrm{act}}
\]
will be used.

However, geometric activity alone is not sufficient to establish
optimality. The conic multiplier must also be considered.

\subsection{SOC Deactivation}

A key research component is determining when the SOC should be
removed from the working geometry.

Unlike a scalar inequality, the SOC has richer boundary structure.
The final implementation will use the conic KKT multiplier and normal
cone conditions to determine whether the cone contributes a genuine
active normal at the candidate solution.

% ============================================================
\section{Degeneracy and Special Cases}
% ============================================================

\subsection{The Point $Lx=0$}

At
\[
Lx=0,
\]
the expression
\[
\frac{Lx}{\|Lx\|_2}
\]
is undefined.

This case must therefore be treated separately.

The SOC boundary at the apex
\[
(t,Lx)=(0,0)
\]
has a different tangent and normal geometry from a smooth nonzero
boundary point.

The implementation will initially handle this case using the exact
SOC membership and normal-cone conditions rather than a tangent
hyperplane.

\subsection{Nearly Active Cones}

Numerical tolerances can lead to oscillation between active and
inactive states.

COSA will therefore use separate activation and deactivation
tolerances, analogous to hysteresis:
\[
\varepsilon_{\mathrm{on}}
<
\varepsilon_{\mathrm{off}}.
\]

\subsection{Degenerate Working Sets}

The working-set KKT matrix may become singular when active constraints
are linearly dependent.

The implementation will detect rank deficiency and investigate:

\begin{itemize}
    \item QR-based rank detection;
    \item null-space methods;
    \item regularization;
    \item dependent-constraint removal.
\end{itemize}

% ============================================================
\section{Development Phases}
% ============================================================

\subsection{Phase I: Polyhedral Baseline}

Implement a conventional active-set solver for
\[
\min_x\;
c^\top x+\frac12x^\top Hx
\]
subject to
\[
Ax\leq b,
\qquad
Ex=d.
\]

Required components:

\begin{itemize}
    \item feasible initialization;
    \item active-set representation;
    \item direction computation;
    \item KKT solve;
    \item multiplier calculation;
    \item constraint addition;
    \item constraint deletion;
    \item ratio test;
    \item termination checks.
\end{itemize}

This implementation will be deliberately simple and will provide the
baseline for debugging COSA.

\subsection{Phase II: SOC Geometry Library}

Implement independent routines for:

\begin{itemize}
    \item SOC membership;
    \item SOC boundary detection;
    \item SOC tangent calculation;
    \item SOC normal calculation;
    \item SOC step feasibility;
    \item exact SOC step interval;
    \item apex handling.
\end{itemize}

These routines should be tested independently using synthetic
directions and analytically known cone intersections.

\subsection{Phase III: COSA Prototype}

Combine the polyhedral active-set method with the SOC geometry.

The prototype should solve the portfolio problem using:

\begin{enumerate}
    \item a working set of linear constraints;
    \item a tangent representation of the SOC;
    \item exact SOC step lengths;
    \item multiplier-based active-set updates.
\end{enumerate}

At this stage, correctness is more important than speed.

\subsection{Phase IV: Full Conic KKT Method}

Replace the preliminary tangent-only treatment with a full
primal-dual conic formulation.

The goal is to obtain a method whose stopping criterion and
working-set logic are derived directly from the SOCP KKT conditions.

\subsection{Phase V: Numerical Linear Algebra}

Optimize the solution of the repeated KKT systems.

Investigate:

\begin{itemize}
    \item sparse LDL$^\top$ factorization;
    \item QR factorization;
    \item null-space methods;
    \item range-space methods;
    \item factorization reuse;
    \item rank-one updates;
    \item scaling;
    \item regularization.
\end{itemize}

\subsection{Phase VI: Warm Starts}

Implement warm starts as a first-class feature.

Given a solution $(x_k,t_k)$ for one problem, initialize the next
problem using:

\begin{itemize}
    \item the previous primal solution;
    \item the previous working set;
    \item previous multipliers where appropriate;
    \item previous KKT factorizations when possible.
\end{itemize}

This is expected to be one of the strongest potential advantages of
COSA.

% ============================================================
\section{Portfolio Test Problems}
% ============================================================

\subsection{Basic Portfolio}

The first model will be
\begin{equation}
\begin{aligned}
\min_{x,t}\quad&
-\mu^\top x+\lambda t\\
\text{s.t.}\quad&
\mathbf{1}^\top x=1,\\
&x\geq0,\\
&\|Lx\|_2\leq t.
\end{aligned}
\tag{8}
\end{equation}

This provides the simplest realistic test.

\subsection{Box Constraints}

Introduce
\[
\ell\leq x\leq u.
\]

These constraints are particularly useful for active-set testing
because many portfolio bounds are expected to become active.

\subsection{Sector Constraints}

Add constraints of the form
\[
A_{\mathrm{sector}}x\leq b_{\mathrm{sector}}.
\]

These create nontrivial combinations of active constraints.

\subsection{Factor Exposure Constraints}

Add
\[
\ell_f\leq Fx\leq u_f.
\]

This provides a useful test of correlated active constraints.

\subsection{Turnover Constraints}

For a previous portfolio $x^{\mathrm{old}}$, turnover can be
represented using auxiliary variables and linear inequalities.

These problems are particularly interesting for warm starts because
successive portfolio rebalancing problems are naturally related.

% ============================================================
\section{Efficient Frontier Experiments}
% ============================================================

One of the principal numerical experiments will solve the sequence

\[
\min_x
\left\{
-\mu^\top x+\lambda\sigma(x)
\right\}
\]

for
\[
\lambda_1,\lambda_2,\ldots,\lambda_N.
\]

The sequence generates points on the risk-return frontier.

For consecutive values
\[
\lambda_k,\lambda_{k+1},
\]
the optimal portfolios and active constraints are expected to be
related.

This allows us to test whether COSA benefits from:

\[
\text{solution}_{k}
\longrightarrow
\text{warm start for solution}_{k+1}.
\]

The main quantities measured will be:

\begin{itemize}
    \item number of active-set iterations;
    \item number of constraints added;
    \item number of constraints removed;
    \item number of KKT factorizations;
    \item total runtime;
    \item KKT residual;
    \item number of iterations saved by warm starts.
\end{itemize}

% ============================================================
\section{Benchmarking}
% ============================================================

\subsection{Reference Solvers}

COSA will be compared with established commercial conic optimization
solvers, initially including MOSEK and Gurobi.

These solvers provide reference solutions for correctness and
performance.

The comparison should distinguish between:

\begin{itemize}
    \item cold-start solves;
    \item warm-start solves;
    \item individual large problems;
    \item sequences of related problems.
\end{itemize}

\subsection{Accuracy Metrics}

For each solution, record:

\begin{itemize}
    \item primal residual;
    \item dual residual;
    \item stationarity residual;
    \item complementarity residual;
    \item objective value;
    \item standard deviation;
    \item expected return.
\end{itemize}

\subsection{Performance Metrics}

Record:

\begin{itemize}
    \item wall-clock time;
    \item number of iterations;
    \item number of KKT solves;
    \item number of active-set changes;
    \item factorization time;
    \item memory usage where relevant.
\end{itemize}

\subsection{Robustness Tests}

Construct instances with:

\begin{itemize}
    \item nearly redundant constraints;
    \item highly correlated assets;
    \item ill-conditioned covariance matrices;
    \item nearly active SOC constraints;
    \item degenerate optimal solutions;
    \item many simultaneously active portfolio bounds.
\end{itemize}

The purpose is to identify failure modes before optimizing performance.

% ============================================================
\section{Numerical Linear Algebra Strategy}
% ============================================================

The active-set approach repeatedly solves systems of the form
\[
\begin{bmatrix}
H&W^\top\\
W&0
\end{bmatrix}
\begin{bmatrix}
d\\
\nu
\end{bmatrix}
=
\begin{bmatrix}
r\\
0
\end{bmatrix}.
\]

Since the working set changes by only a few constraints per iteration,
it may be possible to exploit substantial structure.

\subsection{Initial Implementation}

The first implementation will refactor the KKT matrix at every
iteration.

This minimizes implementation complexity and provides a reliable
reference.

\subsection{Factorization Reuse}

Once the algorithm is correct, investigate updates when:

\begin{itemize}
    \item one constraint is added;
    \item one constraint is removed;
    \item the SOC tangent changes.
\end{itemize}

The SOC tangent changes continuously with $x$, so this part is more
subtle than ordinary linear constraint updates.

\subsection{Scaling}

Scaling will be investigated for:

\begin{itemize}
    \item portfolio variables;
    \item covariance matrices;
    \item linear constraints;
    \item expected returns;
    \item SOC variables.
\end{itemize}

Good scaling is particularly important because the SOC couples $t$
and $Lx$.

% ============================================================
\section{Convergence and Correctness}
% ============================================================

The project will distinguish three levels of validation.

\subsection{Level 1: Feasibility}

Every accepted iterate must satisfy, within tolerance,
\[
Ax\leq b,
\qquad
Ex=d,
\qquad
\|Lx\|_2\leq t.
\]

\subsection{Level 2: Stationarity}

The computed multipliers must satisfy the stationarity equations to
within a prescribed tolerance.

\subsection{Level 3: Conic KKT Optimality}

The final solution must satisfy the full conic KKT conditions.

Because the problem is convex, satisfaction of the KKT conditions
provides a certificate of global optimality under the usual
constraint qualification assumptions.

The final paper should state the precise assumptions required by the
algorithm and distinguish the generic case from degenerate cases.

% ============================================================
\section{Software Architecture}
% ============================================================

The implementation will be modular.

A proposed structure is:

\begin{verbatim}
cosa/
    problem/
        socp.py
        portfolio.py
    geometry/
        soc.py
        tangent.py
        step.py
    active_set/
        working_set.py
        multipliers.py
        updates.py
    linear_algebra/
        kkt.py
        factorization.py
        scaling.py
    solver/
        cosa.py
        initialization.py
        termination.py
    experiments/
        portfolio.py
        frontier.py
        benchmarks.py
    tests/
        test_soc.py
        test_step.py
        test_kkt.py
        test_active_set.py
        test_portfolio.py
\end{verbatim}

The exact language and numerical libraries can be selected at the
start of implementation. Python with NumPy/SciPy is suitable for the
initial prototype; a lower-level implementation can be considered
after the algorithm is stable.

% ============================================================
\section{Testing Strategy}
% ============================================================

\subsection{Unit Tests}

Every geometric primitive should have independent tests.

Examples include:

\begin{itemize}
    \item points inside the SOC;
    \item points outside the SOC;
    \item boundary points;
    \item tangent directions;
    \item directions entering the cone;
    \item directions leaving the cone;
    \item exact step roots.
\end{itemize}

\subsection{Analytical Problems}

Construct small SOCPs whose solutions can be derived analytically.

These problems will be used to validate:

\begin{itemize}
    \item primal solutions;
    \item active sets;
    \item multipliers;
    \item SOC activity;
    \item KKT residuals.
\end{itemize}

\subsection{Cross-Solver Tests}

For every randomly generated test problem, compare COSA against a
reference solver.

The objective values should agree to within the prescribed numerical
tolerance.

% ============================================================
\section{Research Risks}
% ============================================================

\subsection{Risk 1: The SOC Working-Set Concept May Be Insufficient}

The boundary of a second-order cone has richer geometry than a
polyhedral inequality.

If a simple ``SOC active/inactive'' binary state is insufficient, the
algorithm will be generalized to use cone faces and normal-cone
information.

\subsection{Risk 2: Cycling}

As in classical active-set methods, the algorithm may cycle among
working sets.

Potential remedies include:

\begin{itemize}
    \item anti-cycling rules;
    \item multiplier tolerances;
    \item lexicographic selection;
    \item merit-function safeguards.
\end{itemize}

\subsection{Risk 3: Numerical Instability}

Nearly dependent constraints and ill-conditioned covariance matrices
may produce unstable KKT systems.

This will be addressed through scaling, rank detection,
regularization, and robust factorizations.

\subsection{Risk 4: Poor Generic Performance}

COSA may not outperform interior-point SOCP solvers on large,
isolated, generic problems.

This is not considered a failure of the project. The principal
hypothesis is that active-set methods may be particularly attractive
for structured sequences of related problems and problems with a
stable active structure.

% ============================================================
\section{Expected Contributions}
% ============================================================

The project aims to produce four contributions.

\subsection{1. Mathematical Contribution}

A precise formulation of an active-set method for second-order cone
constraints, including the appropriate tangent, normal, and
primal-dual KKT geometry.

\subsection{2. Algorithmic Contribution}

A practical algorithm that combines:

\[
\text{polyhedral working sets}
+
\text{SOC geometry}
+
\text{exact cone-aware step lengths}.
\]

\subsection{3. Software Contribution}

An open and transparent implementation of COSA suitable for research,
experimentation, and extension to other SOCPs.

\subsection{4. Computational Contribution}

A systematic study of when conic active-set methods are competitive
with interior-point methods, with particular emphasis on portfolio
optimization and warm starts.

% ============================================================
\section{Milestones}
% ============================================================

The project will proceed through the following milestones.

\begin{center}
\begin{tabular}{lll}
\toprule
Milestone & Deliverable & Status \\
\midrule
M1 & Mathematical problem specification & Planned \\
M2 & Polyhedral active-set baseline & Planned \\
M3 & SOC geometry module & Planned \\
M4 & Exact SOC step calculation & Planned \\
M5 & First COSA prototype & Planned \\
M6 & Conic KKT formulation & Planned \\
M7 & Robust numerical implementation & Planned \\
M8 & Warm-start capability & Planned \\
M9 & Portfolio benchmark suite & Planned \\
M10 & Comparison with reference solvers & Planned \\
M11 & Final numerical study & Planned \\
M12 & Research paper and software release & Planned \\
\bottomrule
\end{tabular}
\end{center}

% ============================================================
\section{Final Deliverables}
% ============================================================

The project will conclude with the following deliverables.

\begin{enumerate}
    \item A complete mathematical specification of COSA.
    \item A working implementation of the algorithm.
    \item A unit-tested SOC geometry library.
    \item A robust active-set/KKT linear algebra implementation.
    \item A portfolio optimization interface.
    \item A benchmark suite.
    \item Numerical comparisons against established SOCP solvers.
    \item Warm-start experiments along the efficient frontier.
    \item Documentation and reproducible experiments.
    \item A research paper describing the method and results.
\end{enumerate}

% ============================================================
\section{Longer-Term Extensions}
% ============================================================

Once the basic COSA solver is operational, several extensions become
possible.

\subsection{Multiple Second-Order Cones}

The method can be generalized to problems containing several
constraints
\[
(t_j,L_jx)\in\mathcal Q_j,
\qquad
j=1,\ldots,m.
\]

The working set would then track the active geometry of multiple
cones.

\subsection{General SOCPs}

The portfolio structure can be removed entirely, giving a solver for
the generic form
\[
\begin{aligned}
\min_x\quad&
c^\top x\\
\text{s.t.}\quad&
Ax=b,\\
&Gx+h\in\mathcal K,
\end{aligned}
\]
where $\mathcal K$ is a Cartesian product of second-order cones.

\subsection{Mixed Conic Problems}

A later version could investigate combinations of:

\begin{itemize}
    \item linear inequalities;
    \item second-order cones;
    \item rotated second-order cones;
    \item other structured convex cones.
\end{itemize}

\subsection{Large-Scale and Sparse Problems}

For large portfolios and factor models, covariance matrices often
have exploitable low-rank structure. COSA could incorporate
factor-based representations to reduce the cost of the SOC operations.

% ============================================================
\section{Success Criteria}
% ============================================================

The project will be considered successful if COSA satisfies the
following criteria.

\begin{enumerate}
    \item It reliably solves the initial portfolio SOCPs to high
    accuracy.
    
    \item Its termination criterion is based on mathematically
    meaningful conic KKT residuals.
    
    \item Its working-set decisions can be interpreted in terms of
    the active portfolio constraints and SOC geometry.
    
    \item It demonstrates reliable warm starts on sequences of
    related portfolio problems.
    
    \item Its solutions agree with established SOCP solvers.
    
    \item Its numerical behavior is understood on degenerate and
    ill-conditioned problems.
    
    \item The resulting implementation is sufficiently modular to
    support extensions beyond the initial portfolio application.
\end{enumerate}

The most important scientific result need not be that COSA is faster
than every interior-point solver. A valuable result would instead be
a clear characterization of when conic active-set methods work well,
why they work well, and how their geometry differs fundamentally from
polyhedral active-set methods.

% ============================================================
\section{Summary}
% ============================================================

COSA begins with a simple but important modeling observation:
standard deviation is often a more natural risk quantity than
variance because it has the same units as expected return.

The resulting portfolio problem is naturally represented as an SOCP:
\[
\begin{aligned}
\min_{x,t}\quad&
-\mu^\top x+\lambda t\\
\text{s.t.}\quad&
Ax\leq b,\\
&Ex=d,\\
&(t,Lx)\in\mathcal Q.
\end{aligned}
\]

The linear constraints retain the familiar structure of an
active-set method, while the risk term introduces a second-order
cone.

COSA will investigate whether these two structures can be combined
into a single algorithmic framework.

The project therefore has both a practical and a broader algorithmic
objective:

\[
\boxed{
\begin{array}{c}
\text{Portfolio optimization with standard deviation}\\[3pt]
\Downarrow\\[3pt]
\text{Second-order cone formulation}\\[3pt]
\Downarrow\\[3pt]
\text{Conic active-set method}\\[3pt]
\Downarrow\\[3pt]
\text{General active-set methodology for SOCPs}
\end{array}}
\]

The portfolio problem is the test bed. The deeper goal is COSA itself:
a transparent, warm-startable, structure-exploiting active-set
approach to conic optimization.

\end{document}
